% -----------------------------------------------------------------------------
% 3. The Thematic Layer and Its Construction — ressource d'abord, protocole
% ensuite, mesures en dernier. Budget : ~400 mots. Pas de sous-titres numérotés.
% -----------------------------------------------------------------------------
\paragraph{The layer.}
We annotate all 50 English contracts of \corpus{} \cite{lippi2019claudette}:
9,414 sentences, each labelled with one primary theme and zero or more secondary
themes drawn from a 20-theme legal vocabulary (\Tn{20}). Every sentence was
labelled independently by each of three trained
annotators working from a written codebook, and, separately, by four large
language models from three providers (Anthropic, OpenAI and Mistral AI;
Table~\ref{tab:systems}), each run through its vendor's agent interface at default
settings, prompted once with the same vocabulary on the same sentences between
May and August 2026; their outputs were stored and never regenerated. Individual votes are retained rather than collapsed into a
consensus \cite{braun2024differ,plank2022variation,uma2021disagreement}.

\paragraph{Construction.}
Annotators did not post-edit the output of one model. The four judges labelled
every sentence independently, in isolated sessions with no access to one another's
outputs or to any log of the campaign, and a deterministic engine sorted each sentence into
one of five tiers on inter-judge agreement alone, from unanimous to fully split.
Annotators, who could not see one another's work, saw that tier and the competing
candidates but selected from the full 20-theme vocabulary, and frequently did:
34.5\% of gold labels match none of the four proposals, rising to 58.7\% on the
sentences annotators themselves disputed. The
gold standard is built by a three-tier cascade: 46.8\% of sentences are unanimous,
48.3\% carry a majority of at least two of three, and the remaining 4.9\%
(462 sentences) were arbitrated by a human committee --- the fraction of the
corpus on which three trained annotators, given the same codebook, cannot
converge. Crucially,
\emph{model outputs are never parties to the gold}: resolution is strictly
inter-annotator, and the judges are held out as an object of evaluation rather
than a contributor to the reference \cite{thalken2023legal}.

\Tn{14} and \Tn{11} were designed under four criteria --- confusability,
reliability, contractual function and support --- and one constraint fixed in
advance: no merge joins themes whose observed unfairness rates fall in opposite
strata. The merges were reviewed by the legal scholars among the authors;
\Tn{10} breaks that constraint by construction and is reported as a
counterfactual.

\paragraph{Measurement.}
We report Krippendorff's $\alpha$ under two distances
\cite{krippendorff2018content,artstein2008intercoder}: the MASI distance for
set-valued labels \cite{passonneau2006masi}, and the nominal distance on primary
themes alone. The gap between them is what multi-label annotation costs in
reliability \cite{marchal2022multilabel}. Coarser taxonomies (\Tn{14}, \Tn{11},
\Tn{10}) are \emph{projections} of \Tn{20} applied at read time, so every comparison
is paired by construction; confidence intervals come from bootstrap resampling
over documents, never sentences. The merges were designed on 33 contracts and validated on the
17 held out from that design, and the human ceiling is estimated
leave-one-annotator-out against a gold to which the scored annotator did not
contribute \cite{nangia2019human}.
